\section*{الفصل \ref{ch:b2:meiosis-heredity} --- التقسم المنصّف والخلط الوراثي والوراثة}

\begin{solution}{exo:b2:meiosis-heredity:1}
التقسم الخيطي: انقسام واحد، وبلا تزاوج للمتماثلات، وخليتان، كل منهما
مطابقة وراثيًا للأم ($2n \to 2n$). \omterm{def:b2:meiosis-heredity:meiosis}{والتقسم المنصّف}: انقسامان بعد نسخ
واحد، وتتزاوج المتماثلات وتتعابر، وأربع خلايا، فردية الصيغة ومختلفة
وراثيًا كلها ($2n \to n$).
\end{solution}

\begin{solution}{exo:b2:meiosis-heredity:2}
$2^{23} = 8.4\times 10^{6}$؛ و$2^{4} = 16$؛ و$2^{7} = 128$ --- وذلك
قبل أن يضاعف التعابر كلًا منها بلا حد.
\end{solution}

\begin{solution}{exo:b2:meiosis-heredity:3}
$Tt\times Tt$: الأنماط الوراثية $1\,TT : 2\,Tt : 1\,tt$، والأنماط
الظاهرية $3$ طويل : $1$ قزم. و$Tt\times tt$: $1\,Tt : 1\,tt$، أي
$1 : 1$. وهجّن النبات الطويل اختباريًا مع قزم: فإن جاء الخلف كله
طويلًا فهو $TT$، وإن جاء نصفه قزمًا فهو $Tt$ (أو لقّحه ذاتيًا:
فالنمط $TT$ يورّث بصدق، و$Tt$ ينفصل $3 : 1$).
\end{solution}

\begin{solution}{exo:b2:meiosis-heredity:4}
المورّثتان المرتبطتان تقعان على الصبغي نفسه، وتراكيب أليلاتهما
الأبوية ممثَّلة تمثيلًا زائدًا بين الأعراس؛ \omterm{def:b2:meiosis-heredity:linkage}{وتواتر إعادة التركيب} هو
كسر الأعراس المعادة التركيب (أي خلف تهجين اختباري معاد التركيب)؛
\omterm{def:b2:meiosis-heredity:linkage}{والسنتيمورغان} الواحد واحد في المئة من إعادة التركيب. أما المورّثتان
المتباعدتان على صبغي طويل فيقع بينهما تعابر واحد على الأقل في كل
\omterm{def:b2:meiosis-heredity:meiosis}{تقسم منصّف} تقريبًا، وتعشّي التعابرات المتعددة التراكيب، فيبلغ $r$
قيمة $\tfrac12$ --- ولا يتميز ذلك عن \omterm{thm:b2:meiosis-heredity:mixing}{التوزع المستقل}.
\end{solution}

\begin{solution}{exo:b2:meiosis-heredity:5}
المجموع 556؛ والمنتظر $312.75 : 104.25 : 104.25 : 34.75$؛ ومنه $\chi^2 =
2.25^2/312.75 + 3.75^2/104.25 + 3.25^2/104.25 + 2.75^2/34.75 = 0.016
+ 0.135 + 0.101 + 0.218 = 0.47 < 7.81$: أي توافق مع $9 : 3 : 3 :
1$.
\end{solution}

\begin{solution}{exo:b2:meiosis-heredity:6}
في مقابل $1 : 1 : 1 : 1$ (250 لكل صنف): $\chi^2 = (162^2 + 138^2 + 147^2
+ 153^2)/250 = 361 \gg 7.81$: أي مرتبطتان. ومعادات التركيب $103 + 97 = 200$
من 1000: ومنه $r = 0.20$، أي \qty{20}{cM}. وفي حال التنافر يعطي $Ab/aB$ عدد 400
من $Ab$، و400 من $aB$، و100 من $AB$، و100 من $ab$: أي المسافة نفسها،
مع تبادل الصنفين الأبوي والمعاد التركيب.
\end{solution}

\begin{solution}{exo:b2:meiosis-heredity:7}
$d = -\tfrac12\ln(1 - 2r)$: فالقيمة $r = 0.20$ تعطي $-\tfrac12\ln 0.6 =
0.255$، أي \qty{25.5}{cM}؛ و$r = 0.45$ يعطي $-\tfrac12\ln 0.1 = 1.15$،
أي \qty{115}{cM}. وبالمقلوب $r = \tfrac12(1 - e^{-2d})$: فالمسافة \qty{30}{cM} تعطي
$\tfrac12(1 - e^{-0.6}) = 0.226$؛ و\qty{100}{cM} تعطي $\tfrac12(1 -
e^{-2}) = 0.43$. ودون \qty{10}{cM} يكون التصحيح مهملًا؛ وفوق
\qty{30}{cM} يصير لازمًا، وفوق \qty{50}{cM} يتشبع $r$ حتى لا يعود
\omterm{thm:b2:meiosis-heredity:mendel}{تهجين ثنائي} النقاط واحد قادرًا على قياس المسافة --- فتُرسم خريطة
المورّثات المتباعدة بتسلسل المتقاربة.
\end{solution}

\begin{solution}{exo:b2:meiosis-heredity:8}
(أ) $X^wX^w \times X^+Y$: الأبناء كلهم بيض ($X^wY$)، والبنات كلهن
حمر ($X^+X^w$). (ب) $X^+X^w \times X^wY$: نصف الأبناء حمر ونصفهم بيض؛
ونصف البنات حمر ($X^+X^w$) ونصفهن بيض ($X^wX^w$). (ج) البنات
$X^+X^w$ $\times$ إخوتهن $X^wY$: كما في (ب) --- أي نصف كل جنس
أبيض ونصفه أحمر.
\end{solution}

\begin{solution}{exo:b2:meiosis-heredity:9}
إنزيمان على التوالي يصنعان الصبغة؛ ويحتاج اللون الأرجواني إلى \omterm{def:b2:meiosis-heredity:vocabulary}{أليل}
\omterm{def:b2:meiosis-heredity:vocabulary}{سائد} في الموقعين معًا ($C\_P\_$). والسلالتان هما $CCpp$ و$ccPP$
(كل منهما بيضاء لسبب مختلف)؛ والهجين $CcPp$ أرجواني؛
وإذا لُقّح ذاتيًا أعطى $9\,C\_P\_$ أرجوانيًا إلى $7$ أبيض ($3\,C\_pp + 3\,ccP\_ +
1\,ccpp$). والتهجين الاختباري $CcPp\times ccpp$: $1$ أرجواني ($CcPp$) : $3$
أبيض.
\end{solution}

\begin{solution}{exo:b2:meiosis-heredity:10}
الأبويان $+++$ و$abc$؛ والتعابران المضاعفان $+b+$ و$a+c$؛ والمورّثة
التي تبدلت هي $b$: فالترتيب $a$--$b$--$c$. ومنه $r_{ab} = (11 + 9 + 2)/1202
= 0.018$ (أي \qty{1.8}{cM})؛ و$r_{bc} = (118 + 112 + 2)/1202 = 0.193$
(أي \qty{19.3}{cM}). والمضاعفان المنتظران $0.018\times 0.193\times 1202 =
4.2$؛ والمرصودان 2: فالتصادف $0.47$، والتداخل $0.53$.
\end{solution}

\begin{solution}{exo:b2:meiosis-heredity:11}
أكياس الانفصال الثاني $300/1000 = \qty{30}{\%}$؛ والمسافة $=
\tfrac12\times 30 = \qty{15}{cM}$. فالتعابر بين المورّثة والقسيم
المركزي لا يشمل إلا اثنتين من الكروماتيدات الأربع، فالكيس الذي
يُظهر انفصالًا في الانقسام الثاني يحمل كروماتيدتين معادتي التركيب
واثنتين أبويتين: \omterm{def:b2:meiosis-heredity:linkage}{فتواتر إعادة التركيب} نصف تواتر تلك الأكياس. وأما
$2 : 4 : 2$: فالتعابر بين الموقع والقسيم المركزي يبادل الكروماتيدتين
الداخليتين، فيحمل كل نصف مغزل في \omterm{def:b2:meiosis-heredity:meiosis}{التقسم المنصّف} الأول كروماتيدة
$A$ وأخرى $a$؛ ثم يفصلهما \omterm{def:b2:meiosis-heredity:meiosis}{التقسم المنصّف} الثاني بالترتيب $A\,a \mid
a\,A$ (أو العكس)، ويضاعف التقسم الخيطي الأخير كل بوغ:
$2 : 4 : 2$.
\end{solution}

\begin{solution}{exo:b2:meiosis-heredity:12}
أمها حاملة إلزامًا (إذ تلقّت الصبغي X الوحيد لأبيها)، فتكون المرأة
قد تلقت \omterm{def:b2:meiosis-heredity:vocabulary}{الأليل} باحتمال $\tfrac12$.
وابنان سليمان: فللحاملة احتمال $(\tfrac12)^2 = \tfrac14$
في ابنين سليمين، ولغير الحاملة احتمال 1. وبقاعدة بايز:
$P(\text{حاملة}\mid\text{المعطيات}) = (\tfrac12\times\tfrac14)/(\tfrac12
\times\tfrac14 + \tfrac12\times 1) = \tfrac{1/8}{5/8} = \tfrac15$.
\end{solution}

\begin{solution}{pb:b2:meiosis-heredity:1}
\textbf{1.} الأبوان $vg/vg\;+/+$ و$+/+\;e/e$؛ والجيل $F_1$ هو $+/vg\;+/e$؛
والأعراس $+\,+$ و$+\,e$ و$vg\,+$ و$vg\,e$، كل منها $\tfrac14$.
\textbf{2.} $900 : 300 : 300 : 100$.
\textbf{3.} $\chi^2 = 8^2/900 + 9^2/300 + 4^2/300 + 3^2/100 = 0.07 +
0.27 + 0.05 + 0.09 = 0.48 < 7.81$: أي إن \omterm{thm:b2:meiosis-heredity:mixing}{التوزع المستقل} قائم.
\textbf{4.} \omterm{def:b2:meiosis-heredity:vocabulary}{المتماثل اللواقح} البري في الموقعين هو $\tfrac{1}{16}$ من
الجميع، والبرية $\tfrac{9}{16}$: فالنسبة $\tfrac19$.
\textbf{5.} برية، وأثرية، وأبنوسية، وأثرية أبنوسية، $\tfrac14$ لكل صنف.
\textbf{6.} \omterm{def:b2:meiosis-heredity:vocabulary}{النمط الظاهري} جداء \omterm{def:b2:meiosis-heredity:vocabulary}{النمط الوراثي} في المحيط (فالحرارة
تؤثر في إنزيمات الصبغة): \omterm{def:b2:meiosis-heredity:vocabulary}{فالنمط الوراثي} يحدد معيار استجابة لا صفة
ثابتة.
\textbf{7.} $b\,+/+\,vg$ (تنافر): إذ جاء $b$ مع $+_{vg}$ من أحد
الأبوين، وجاء $vg$ مع $+_b$ من الآخر.
\textbf{8.} الأبوية: السوداء ($b\,+$) 965 والأثرية ($+\,vg$) 944؛
والمعادة التركيب: السوداء الأثرية ($b\,vg$) 206 والبرية ($+\,+$) 185.
فالصبغي البري $+\,+$ لم يكن موجودًا في الجيل $F_1$؛ ولا يمكن
صنعه إلا بتعابر.
\textbf{9.} $r = 391/2300 = 0.17$: أي \qty{17}{cM}.
\textbf{10.} $d = -\tfrac12\ln(1 - 0.34) = -\tfrac12\ln 0.66 = 0.21$:
أي \qty{21}{cM}.
\textbf{11.} لا تعابر عند ذكور \emph{Drosophila}:
فأعراس الذكر لا تحمل إلا الصبغيين الأبويين.
\textbf{12.} الجيل $F_1$ هو $b\,vg/+\,+$ (تزاوج): فالصنفان الأبويان
بريان وأسودان أثريان، \qty{41.5}{\%} لكل منهما؛ والمعادان التركيب
أسود وأثري، \qty{8.5}{\%} لكل منهما.
\textbf{13.} $X^hX^+$: فالابنة تتلقى الصبغي X الوحيد لأبيها، وهو
يحمل $h$، وتتلقى صبغي X سليمًا من أمها غير المصابة.
\textbf{14.} $\tfrac12$ (أن يكون ابنًا) $\times$ $\tfrac12$ (أن يتلقى $X^h$) $=
\tfrac14$.
\textbf{15.} $\tfrac12$: إذ تتلقى أحد صبغيي X الأموميين؛ ويعطيها
الأب صبغي X سليمًا.
\textbf{16.} كل طفل غير مصاب باحتمال $\tfrac34$:
ومنه $(\tfrac34)^3 = \tfrac{27}{64} = 0.42$.
\textbf{17.} الابن يتلقى Y أبيه لا X. والمرأة لا تصاب إلا بالنمط
$X^hX^h$: أي أب مصاب وأم حاملة (أو مصابة)، وهذا نادر.
\textbf{18.} الاحتمال القبلي $\tfrac12$. واحتمال ابنين سليمين
$\tfrac14$ إن كانت حاملة، و1 إن لم تكن: فالاحتمال البعدي $(\tfrac12\times
\tfrac14)/(\tfrac12\times\tfrac14 + \tfrac12) = \tfrac15$.
\textbf{19.} الصبغي X الوحيد للابن يأتي من أمه، والصبغي Y لا يحمل
شيئًا من هذه المورّثات، فيُظهر كل ابن عروسًا أمومية واحدة بلا قناع:
فالتهجين تهجين اختباري لنفسه.
\textbf{20.} الأبويان $+++$ (853) و$ywm$ (833)؛ والتعابران المضاعفان
$+w+$ (2) و$y+m$ (2)؛ والمورّثة التي تبدلت هي $w$: فالترتيب $y$--$w$--$m$.
\textbf{21.} $r_{yw} = (24 + 26 + 2 + 2)/2000 = 0.027$؛ و$r_{wm} = (128
+ 132 + 2 + 2)/2000 = 0.132$. والخريطة: $y$ --- \qty{2.7}{cM} --- $w$ ---
\qty{13.2}{cM} --- $m$.
\textbf{22.} المنتظر $0.027\times 0.132\times 2000 = 7.1$؛ والمرصود
4: فالتصادف $0.56$، والتداخل $0.44$.
\textbf{23.} معادات التركيب $y$--$m$ هي $y++$ و$+wm$ و$yw+$ و$++m$: أي $(24 +
26 + 128 + 132)/2000 = 0.155$، في مقابل $0.027 + 0.132 = 0.159$:
فالتعابرات المضاعفة الأربعة تعيد التركيب الأبوي $y$--$m$ ولا
تُحصى، أي $2\times 4/2000 = 0.004$.
\textbf{24.} تتلقى البنات صبغي $y\,w\,m$ الأبوي وصبغيًا أموميًا
واحدًا: فكلهن متغايرات اللواقح وبريات في \omterm{def:b2:meiosis-heredity:vocabulary}{النمط الظاهري}، مهما كان
الصبغي المعاد التركيب الذي يحملنه؛ فإعادة التركيب غير مرئية.
\textbf{25.} $y$ --- \qty{2.7}{cM} --- $w$ --- \qty{13.2}{cM} --- $m$؛
ومعامل التصادف $0.56$، والتداخل $0.44$.
\end{solution}
